\expandafter\gdef\csname SNside1\endcsname{\node[anchor=north west,inner sep=0pt,text width=48mm,font=\SideFont,align=left] at ([xshift=138mm,yshift=-466.85pt]current page.north west) {\textcolor{SNink}{\textbf{最小反例 · 数量相同仍不合法}}\par\vspace{1.5mm}顺序输入a、b、c，要先输出c就必须把三者全部入栈；弹出c后栈顶为b，接着要求a便失败。输出c、a、b与输入含相同元素，却违反栈顶约束。\par\smallskip\begin{tikzpicture}[x=1mm,y=1mm,>=Stealth,node distance=5mm]
\node[draw=SNink!50,fill=SNblue,rounded corners=1mm,text width=34mm,align=center,inner sep=1.2mm,font=\fontsize{7.5}{9}\selectfont] (m0) at (0,0) {底a};
\node[draw=SNink!50,fill=SNblue,rounded corners=1mm,text width=34mm,align=center,inner sep=1.2mm,font=\fontsize{7.5}{9}\selectfont] (m1) at (0,-11) {顶b};
\node[draw=SNink!50,fill=SNblue,rounded corners=1mm,text width=34mm,align=center,inner sep=1.2mm,font=\fontsize{7.5}{9}\selectfont] (m2) at (0,-22) {要求a};
\draw[->,SNink] (m0) -- (m1);
\draw[->,SNink] (m1) -- (m2);
\end{tikzpicture}
\par{\fontsize{7.5}{9}\selectfont c已弹出，b挡在a上方}};}
\expandafter\gdef\csname SNside2\endcsname{\node[anchor=north west,inner sep=0pt,text width=48mm,font=\SideFont,align=left] at ([xshift=138mm,yshift=-428.92pt]current page.north west) {\textcolor{SNink}{\textbf{错项排除 · 共享栈只在相邻时满}}\par\vspace{1.5mm}两栈从数组两端向中间增长，未用单元恰位于两栈顶之间。两栈顶相邻已经没有空位；等到相等才判满，就已让两个栈占用同一个位置。};}
\expandafter\gdef\csname SNside3\endcsname{\node[anchor=north west,inner sep=0pt,text width=48mm,font=\SideFont,align=left] at ([xshift=138mm,yshift=-95.00pt]current page.north west) {\textcolor{SNink}{\textbf{错项排除 · 双端队列要看受限端}}\par\vspace{1.5mm}输入受限和输出受限双端队列是两种不同约束。普通栈的合法输出在一般双端队列中能实现，但不能倒推出任意指定输出都能实现；仍要跟踪两端状态。};}
\expandafter\gdef\csname SNside4\endcsname{\node[anchor=north west,inner sep=0pt,text width=48mm,font=\SideFont,align=left] at ([xshift=138mm,yshift=-95.00pt]current page.north west) {\textcolor{SNink}{\textbf{不变量 · 尾指针不能悬空}}\par\vspace{1.5mm}带头结点链队列始终让rear指向最后一个数据结点，空时指向头结点。删去唯一数据结点后，rear必须回到front；否则下次入队会经已释放地址写内存。};}
\expandafter\gdef\csname SNside5\endcsname{\node[anchor=north west,inner sep=0pt,text width=48mm,font=\SideFont,align=left] at ([xshift=138mm,yshift=-95.00pt]current page.north west) {\textcolor{SNink}{\textbf{课程连接 · 偏移先统一地址单位}}\par\vspace{1.5mm}数组公式先数元素，再乘每元素宽度。计组中的字节编址与按字编址会改变地址单位；同样的元素偏移，在两种编址方式下不能直接共用一个数值。};}
\expandafter\gdef\csname SNside6\endcsname{\node[anchor=north west,inner sep=0pt,text width=48mm,font=\SideFont,align=left] at ([xshift=138mm,yshift=-95.00pt]current page.north west) {\textcolor{SNink}{\textbf{错项排除 · 三对角首尾行少一个}}\par\vspace{1.5mm}n至少为2时，首行和末行各只保存2个位置，中间行才各3个。若把第i行之前全部按3个计，会从第一行就多算；先列2+3(i\ensuremath{-}2)再加本行偏移。};}
\expandafter\gdef\csname SNside7\endcsname{\node[anchor=north west,inner sep=0pt,text width=48mm,font=\SideFont,align=left] at ([xshift=138mm,yshift=-95.00pt]current page.north west) {\textcolor{SNink}{\textbf{最小反例 · 括号计数看不到嵌套}}\par\vspace{1.5mm}({[}){]}中每一类左右括号数都相等，但第三个字符要求圆括号配对，栈顶却是方括号。配对必须遵守最近未完成对象先结束，不能在扫描后任意重排。};}
\expandafter\gdef\csname SNside8\endcsname{\node[anchor=north west,inner sep=0pt,text width=48mm,font=\SideFont,align=left] at ([xshift=138mm,yshift=-95.00pt]current page.north west) {\textcolor{SNink}{\textbf{不变量 · top含义决定增减位置}}\par\vspace{1.5mm}top指栈顶时，空状态为\ensuremath{-}1，入栈先加再写；top指下一空位时，空状态为0，入栈先写再加。用空栈和单元素两幅状态检查，比只背push语句更能应对改约定题。};}
\expandafter\gdef\csname SNside9\endcsname{\node[anchor=north west,inner sep=0pt,text width=48mm,font=\SideFont,align=left] at ([xshift=138mm,yshift=-95.00pt]current page.north west) {\textcolor{SNink}{\textbf{不变量 · 先换到已保存半边}}\par\vspace{1.5mm}对称矩阵访问未存的半边时，先交换i、j，再按保存方向计数。坐标转换解决的是表示对应；数组从0还是1起始解决的是编号，两步不能互相替代。};}
\expandafter\gdef\csname SNside10\endcsname{\node[anchor=north west,inner sep=0pt,text width=48mm,font=\SideFont,align=left] at ([xshift=138mm,yshift=-95.00pt]current page.north west) {\textcolor{SNink}{\textbf{错项排除 · FIFO只是具体策略}}\par\vspace{1.5mm}队列ADT规定先进先出，操作系统的``就绪队列''名称却不保证所有调度算法都按FIFO选择。先来先服务可借FIFO说明，优先级等策略需要额外排序或选择规则。};}
\expandafter\gdef\csname SNside11\endcsname{\node[anchor=north west,inner sep=0pt,text width=48mm,font=\SideFont,align=left] at ([xshift=138mm,yshift=-95.00pt]current page.north west) {\textcolor{SNink}{\textbf{最小反例 · 同一对下标可能空也可能满}}\par\vspace{1.5mm}数组长度2、front=rear=0可表示刚初始化的空队列，也可在允许全用两个单元时表示已入队2个元素。只有两个模下标丢失了一圈的信息，必须牺牲单元或增加状态。};}
\AddToShipoutPictureFG{\begin{tikzpicture}[remember picture,overlay]\ifcsname SNside\thepage\endcsname\csname SNside\thepage\endcsname\fi\end{tikzpicture}}